%==============================================================================
%  INTRODUCTION GÉNÉRALE — non numérotée (avant le chapitre 1)
%==============================================================================
\chapter*{\textbf{Introduction générale}}
\addcontentsline{toc}{chapter}{Introduction générale}
\mtcaddchapter % aligne le compteur minitoc (pas de Sommaire sur l'intro)
\label{chap:introduction}
\markboth{Introduction générale}{Introduction générale}
\thispagestyle{plain}

L'évolution rapide des technologies numériques et le développement des
plateformes éducatives en ligne ont profondément transformé les pratiques
d'apprentissage, en particulier à travers la diffusion des cours en ligne
ouverts et massifs (MOOC, \emph{Massive Open Online Courses}). Ces dispositifs
permettent aujourd'hui de collecter, diffuser et personnaliser des contenus
pédagogiques à grande échelle, offrant à des millions d'apprenants un accès
élargi à la formation, indépendamment des contraintes géographiques
\cite{breslow2013studying,daniel2012making}. La crise sanitaire de 2020 a
encore accéléré cette dynamique, faisant de l'apprentissage en ligne une
composante désormais structurelle des systèmes éducatifs contemporains.

Toutefois, cette évolution s'accompagne d'un constat largement documenté :
les MOOC souffrent de taux d'abandon particulièrement élevés. Une part
importante des inscrits interrompt son parcours dès les premières semaines,
ce qui limite fortement l'efficacité pédagogique réelle de ces dispositifs
\cite{onah2014dropout,jordan2015massive}. Parmi les facteurs explicatifs,
la baisse de l'\textbf{engagement} de l'apprenant occupe une place centrale
\cite{kizilcec2013deconstructing,hew2016promoting}. Dans une salle de classe
traditionnelle, l'enseignant perçoit en permanence des signaux fins (regard,
posture, expression) et adapte spontanément son discours. Cette boucle de
régulation disparaît presque entièrement dans les environnements en ligne de
masse, où l'apprenant fait face à un contenu figé, identique pour tous et
insensible à son état du moment.

Dans ce contexte, plusieurs travaux en informatique affective et en
psychologie de l'apprentissage ont établi que les émotions ne sont pas
périphériques à la cognition, mais en constituent un rouage essentiel
\cite{picard1997affective,pekrun2006control}. Des états tels que l'ennui, la
confusion ou la frustration influencent directement la profondeur du traitement
de l'information, la persévérance et la mémorisation
\cite{dmello2012dynamics,baker2010better}. Une détection inefficace de ces
états, combinée à l'absence d'adaptation pédagogique, peut entraîner une
dégradation de l'expérience d'apprentissage et favoriser le décrochage. Cette
problématique est d'autant plus critique que les décisions pédagogiques doivent
idéalement intervenir en temps utile, avant que le désengagement ne s'installe.

Ainsi, ce mémoire s'inscrit dans cette problématique en proposant une approche
visant à réunir, au sein d'une même chaîne de traitement, la perception de
l'engagement émotionnel et la génération adaptative de contenu pédagogique.

Face à ces enjeux, la question centrale de ce travail est la suivante :

\begin{center}
\textit{Comment concevoir un dispositif capable d'inférer automatiquement
l'état d'engagement d'un apprenant à partir de signaux multimodaux, puis de
générer en réponse un contenu pédagogique adapté à cet état, tout en
garantissant la cohérence et la traçabilité de la décision pédagogique
sous-jacente~?}
\end{center}

Cette question se décompose en quatre sous-questions opérationnelles :

\begin{itemize}
  \item Comment estimer de manière fiable l'état d'engagement à partir de
  signaux multimodaux hétérogènes (expressions faciales issues de
  \textsc{DAiSEE}, signal vocal issu de \textsc{IEMOCAP}), et quel gain la
  fusion apporte-t-elle par rapport à une modalité unique~?

  \item Comment formaliser une correspondance pédagogiquement justifiable
  entre les états d'engagement détectés et les stratégies d'intervention,
  c'est-à-dire un modèle \emph{engagement\,$\rightarrow$\,adaptation}~?

  \item Dans quelle mesure un grand modèle de langue, guidé par des patrons
  de prompts structurés et ancré par un mécanisme de récupération augmentée
  (RAG), peut-il produire un contenu pédagogique pertinent et adapté à l'état
  d'engagement visé~?

  \item Selon quel protocole peut-on évaluer, de façon honnête et
  multi-niveaux, la qualité de la perception, la pertinence de la décision et
  la qualité du contenu généré~?
\end{itemize}

L'analyse de la littérature montre que, malgré les progrès importants réalisés
dans les domaines de la détection de l'engagement et de l'IA générative
appliquée à l'éducation, ces deux axes restent majoritairement traités de
manière indépendante. D'une part, les modèles de perception multimodale
permettent d'inférer l'état affectif d'un apprenant avec une fiabilité
croissante \cite{dmello2015review,whitehill2014faces,gupta2016daisee}.
D'autre part, les grands modèles de langue ouvrent la voie à une génération
de contenu pédagogique personnalisée à grande échelle
\cite{brown2020language,kasneci2023chatgpt}. Cependant, ces approches
demeurent souvent séquentielles ou fragmentées~: la détection n'est pas
systématiquement reliée à une décision pédagogique explicite, tandis que la
génération n'est pas toujours conditionnée par un état d'engagement mesuré,
ce qui limite la capacité des systèmes à restaurer une forme de régulation
pédagogique automatique.

Pour répondre à ces limitations, ce mémoire propose un framework hybride nommé
\emolearn{}, combinant trois briques complémentaires. La couche de
\textbf{perception} estime l'état d'engagement à partir de signaux visage et
voix, en s'appuyant sur les corpus \textsc{DAiSEE} et \textsc{IEMOCAP} et sur
des modèles \emph{fine-tunés} vers la taxonomie des quatre états
d'engagement~: un Vision Transformer (ViT) pour le visage et un modèle
wav2vec2 pour la voix, combinés par un moteur de fusion multimodale. La couche
de \textbf{décision} formalise un modèle
\emph{engagement\,$\rightarrow$\,adaptation} reliant chaque état à un type de
contenu et à un patron de génération. Enfin, la couche de \textbf{génération}
produit un contenu pédagogique adaptatif au moyen d'un grand modèle de langue,
guidé par une bibliothèque de vingt patrons structurés et ancré par un module
RAG sur une base documentaire pédagogique. Ainsi, \emolearn{} vise à passer
d'un contenu figé à une adaptation réellement conditionnée par l'état
d'engagement détecté.

En résumé, le module de perception anticipe l'état émotionnel de l'apprenant,
tandis que les modules de décision et de génération exploitent cette estimation
pour produire une intervention pédagogique ciblée, traçable et reproductible.

Les principales contributions de ce travail sont les suivantes :

\begin{enumerate}
  \item Proposition d'une architecture multimodale de bout en bout reliant
  perception, décision et génération au sein d'une même chaîne.

  \item Formalisation d'un modèle \emph{engagement\,$\rightarrow$\,adaptation}
  assorti d'une bibliothèque de vingt patrons de génération structurés.

  \item Développement d'un module RAG pédagogique ancrant la génération sur
  une base documentaire vectorisée.

  \item Fine-tuning supervisé des modèles de perception (ViT pour le visage et
  wav2vec2 pour la voix) vers les quatre états d'engagement \textsc{DAiSEE}.

  \item Évaluation honnête et instrumentée de la fusion multimodale, incluant
  l'analyse des compromis entre exactitude globale et sensibilité aux états
  critiques.

  \item Prototype logiciel complet et reproductible, accompagné d'un protocole
  d'évaluation multi-niveaux.
\end{enumerate}

%==============================================================================
\section*{Méthodologie générale}
\label{sec:methodologie}

La démarche adoptée relève de la \emph{recherche orientée conception}
(\emph{design science research}), paradigme particulièrement adapté à la
conception d'artefacts informatiques répondant à un problème identifié
\cite{hevner2004design}. Elle articule l'analyse d'un problème réel (le
décrochage lié à la baisse d'engagement dans les MOOC), la construction
d'un artefact (le framework \emolearn{}) et son évaluation au regard
d'objectifs explicites.

Concrètement, le travail s'organise autour d'une chaîne de traitement en trois
étages : un étage de \emph{perception} (FER, SER, fusion), un étage de
\emph{décision} (règles temporelles et moteur pédagogique) et un étage de
\emph{génération} (RAG et LLM).

Le projet s'appuie sur deux corpus publics de référence. \textsc{DAiSEE}
\cite{gupta2016daisee} constitue le socle engagement (taxonomie et scénarios),
tandis qu'\textsc{IEMOCAP} \cite{busso2008iemocap} apporte la modalité vocale
et, via un sous-corpus audiovisuel aligné, le support de validation de la
fusion. Le prototype repose sur une architecture client-serveur, des modèles
de perception fine-tunés (ViT et wav2vec2) vers les quatre états
\textsc{DAiSEE}, un LLM accessible par API et un module RAG (ChromaDB).

%==============================================================================
\section*{Structure du mémoire}
\label{sec:structure}

Le présent mémoire s'ouvre par la présente \textbf{introduction générale},
puis s'organise en six chapitres numérotés.

Le \textbf{chapitre~1} pose les fondements conceptuels de l'apprentissage
adaptatif et de l'engagement (MOOC, émotions, détection multimodale, IA
générative).

Le \textbf{chapitre~2} propose un état de l'art critique et positionne
\emolearn{} par rapport aux travaux existants.

Le \textbf{chapitre~3} est consacré à la modélisation et à la conception du
framework : architecture globale, modèle
\emph{engagement\,$\rightarrow$\,adaptation}, génération fondée sur des patrons
et un module RAG, ainsi que conception du prototype.

Le \textbf{chapitre~4} décrit l'implémentation technique de bout en bout.

Le \textbf{chapitre~5} expose la validation et les résultats selon un protocole
multi-niveaux.

Le \textbf{chapitre~6}, enfin, synthétise les contributions, expose les limites
et trace des perspectives.

Le mémoire se termine par la bibliographie.
